% =====================================================================
%  File: sections/3_4_kien_truc_de_xuat.tex
%  Mục đích: Phần 3.4 - Kiến trúc và cách giải quyết đề xuất (file cha)
% =====================================================================

\subsection{Kiến trúc và cách giải quyết đề xuất}
\label{sec:kien-truc-de-xuat}

Sau khi phân tích các hệ thống liên quan ở phần~\ref{sec:khao-sat-he-thong}, em chọn hướng tiếp cận \textbf{evolutionary redesign} cho prototype Raptix: kế thừa Decepticon làm nền tảng vì kiến trúc multi-agent theo kill chain và sandbox biệt lập của Decepticon đã được kiểm chứng~\cite{decepticon}, đồng thời tiếp thu các bài học quan trọng từ Incalmo~\cite{incalmo} về tách planning/execution và quản lý state ngoài context. Cách tiếp cận này có rủi ro thấp hơn so với việc viết lại từ đầu, và cho phép sử dụng benchmark làm hàng rào non-regression trong suốt quá trình phát triển.

Đóng góp thiết kế của Raptix tập trung vào hai điểm khác biệt cốt lõi so với Decepticon gốc:

\begin{itemize}
    \item \textbf{Attack Knowledge Graph (AKG):} thay thế cách tổ chức knowledge graph hiện tại của Decepticon bằng một đồ thị \emph{hợp nhất} chứa tất cả các thực thể (asset, identity, credential, vulnerability, finding, action, evidence) và quan hệ giữa chúng, đóng vai trò nguồn sự thật duy nhất cho mọi tác tử.
    \item \textbf{Mô hình khai thác hai pha:} phân tách quy trình khai thác thành pha Initial Foothold (khai thác theo chiều ngang, breadth-first) và pha Network Propagation (khai thác theo chiều sâu, depth-first), với cơ chế chuyển pha tường minh dựa trên trạng thái của AKG.
\end{itemize}

Ba phần con của mục này lần lượt trình bày: kiến trúc tổng thể Raptix (\ref{sec:kien-truc-tong-the}), thiết kế chi tiết AKG (\ref{sec:attack-knowledge-graph}) và mô hình khai thác hai pha (\ref{sec:mo-hinh-hai-pha}).

% =====================================================================
%  File: sections/3_4_1_kien_truc_tong_the.tex
%  Mục đích: Phần 3.4.1 - Kiến trúc tổng thể Raptix
% =====================================================================

\subsubsection{Kiến trúc tổng thể Raptix}
\label{sec:kien-truc-tong-the}

Raptix được tổ chức thành năm lớp chức năng có ranh giới rõ ràng, giao tiếp với nhau qua các interface được định nghĩa tường minh. Việc tách lớp này giúp hệ thống dễ kiểm thử, dễ mở rộng và cho phép thay thế từng thành phần (ví dụ: đổi LLM provider, đổi graph backend) mà không ảnh hưởng đến các lớp khác.

\paragraph{Năm lớp chức năng}

\begin{itemize}
    \item \textbf{Lớp Mission và Policy:} chịu trách nhiệm nhận mission từ operator, kiểm tra scope và policy, nạp thông tin ban đầu vào AKG, và duy trì ranh giới mục tiêu trong suốt phiên làm việc. Đây là lớp bảo vệ đầu tiên, đảm bảo mọi hành động tiếp theo đều nằm trong phạm vi được cấp quyền.
    \item \textbf{Lớp Planner và Orchestrator:} thực hiện phân rã nhiệm vụ, chọn chiến lược khai thác theo pha hiện tại, và điều phối các sub-agent chuyên biệt. Planner chỉ thao tác ở mức trừu tượng cao — chọn \emph{làm gì} — chứ không sinh ra lệnh shell cấp thấp.
    \item \textbf{Lớp Specialist Agent:} gồm các sub-agent chuyên biệt theo domain (recon, web, AD, cloud, post-exploit...). Mỗi agent chỉ được khởi tạo khi cần, với context window sạch và nhiệm vụ được định nghĩa rõ ràng. Agent \emph{không} giữ state riêng ngoài context tạm thời — mọi state đều đọc/ghi qua AKG.
    \item \textbf{Lớp Tool Executor và Sandbox:} chịu trách nhiệm thực thi lệnh trong sandbox Kali Linux biệt lập. Mọi tool call phải đi qua scope guard và audit logger trước khi được đưa xuống sandbox. Tầng này hỗ trợ các phiên tương tác dài thông qua tmux sessions, cho phép làm việc với các công cụ như \texttt{msfconsole} hay \texttt{sliver-client}.
    \item \textbf{Lớp Knowledge và Verification:} bao gồm AKG (lưu trữ tri thức tấn công), event stream (audit và replay), verifier (xác minh finding độc lập), và reporting engine (sinh báo cáo cuối cùng dựa trên bằng chứng).
\end{itemize}

\paragraph{Sơ đồ luồng xử lý}

\begin{figure}[H]
\centering
\begin{tikzpicture}[
    node distance=8mm and 12mm,
    box/.style={draw, rounded corners, align=center, minimum height=9mm, inner sep=4pt, font=\small},
    graph/.style={draw, rounded corners, fill=blue!10, align=center, minimum height=12mm, inner sep=5pt, font=\small\bfseries},
    arrow/.style={-{Latex[length=2mm]}, thick}
]
\node[box, fill=gray!15] (mission) {Mission\\+ Scope\\+ Policy};
\node[box, right=of mission] (planner) {Planner /\\Orchestrator};
\node[box, right=of planner] (agent) {Specialist\\Agent};
\node[box, right=of agent] (executor) {Tool Executor\\+ Sandbox};
\node[graph, below=14mm of executor] (akg) {Attack Knowledge Graph\\(assets, identities, credentials,\\findings, actions, evidence)};
\node[box, left=of akg] (verifier) {Verifier\\+ Reporting};
\node[box, left=of verifier] (bench) {Benchmark\\Harness};

\draw[arrow] (mission) -- (planner);
\draw[arrow] (planner) -- (agent);
\draw[arrow] (agent) -- (executor);
\draw[arrow] (executor) -- (akg);
\draw[arrow] (akg) -- (verifier);
\draw[arrow] (verifier) -- (bench);
\draw[arrow] (akg.west) -- ++(-5mm,0) |- (planner.south);
\end{tikzpicture}
\caption{Luồng xử lý logic của Raptix với AKG ở trung tâm}
\label{fig:raptix-architecture}
\end{figure}

\paragraph{Nguyên tắc thiết kế}

Sáu nguyên tắc thiết kế định hình toàn bộ kiến trúc Raptix:

\begin{itemize}
    \item \textbf{Scope là dữ liệu hạng nhất:} mọi tool call phải được kiểm tra theo scope và policy trước khi chạy. Scope không phải là ràng buộc ngầm trong prompt mà là dữ liệu có cấu trúc được lưu trong AKG và được enforce ở tầng executor.
    \item \textbf{LLM context là cache làm việc, không phải nguồn sự thật:} nguồn sự thật duy nhất là AKG. Điều này cho phép agent khởi động lại, thay đổi model, hoặc chạy song song mà vẫn giữ tính nhất quán của state.
    \item \textbf{Tool execution có định danh và vòng đời rõ ràng:} mỗi tool call có execution\_id, trạng thái (pending, running, completed, failed, cancelled), timeout và output cap, cho phép mission dài có thể resume sau sự cố.
    \item \textbf{Agent specialization chỉ dùng khi có lợi ích rõ:} multi-agent chỉ được kích hoạt khi việc phân tách nhiệm vụ giảm complexity hoặc tăng parallelism đủ để bù đắp coordination overhead. Không mặc định multi-agent cho mọi task.
    \item \textbf{Finding có lifecycle chặt chẽ:} mỗi finding đi qua các trạng thái candidate $\rightarrow$ validated $\rightarrow$ verified $\rightarrow$ reported. Không sinh finding ``chắc chắn'' chỉ từ suy luận của LLM.
    \item \textbf{Graph là nguồn sự thật duy nhất:} mọi agent đọc/ghi qua AKG; không agent nào giữ state riêng ngoài context tạm thời. Điều này đảm bảo tính nhất quán khi nhiều agent chạy song song.
\end{itemize}


% =====================================================================
%  File: sections/3_4_2_attack_knowledge_graph.tex
%  Mục đích: Phần 3.4.2 - Attack Knowledge Graph (AKG)
% =====================================================================

\subsubsection{Attack Knowledge Graph (AKG)}
\label{sec:attack-knowledge-graph}

Attack Knowledge Graph (AKG) là đóng góp thiết kế cốt lõi của Raptix. Đây là một đồ thị có hướng, có nhãn, được thiết kế để lưu trữ \emph{tất cả} các thực thể và quan hệ mà một red team agent cần trong suốt mission, thay thế cho mô hình state phẳng hoặc nhồi toàn bộ vào context LLM.

\paragraph{Các loại node trong AKG}

AKG phân loại node theo bảy nhóm chính, tương ứng với các thực thể mà một red teamer thao tác trong thực tế:

\begin{table}[H]
\centering
\caption{Các loại node trong Attack Knowledge Graph}
\label{tab:akg-nodes}
\small
\renewcommand{\arraystretch}{1.3}
\begin{tabular}{|p{2.6cm}|p{4.5cm}|p{8.5cm}|}
\hline
\textbf{Nhóm node} & \textbf{Ví dụ} & \textbf{Vai trò} \\
\hline
Asset & Host, Port, Service, Endpoint, Application & Mô tả bề mặt tấn công đã quan sát được. \\
\hline
Identity & User, Account, Group, Role & Đối tượng xác thực và phân quyền. \\
\hline
Credential & Password, NTLM hash, Kerberos ticket, SSH key, Session cookie & Vật liệu xác thực thu được trong quá trình khai thác. \\
\hline
Vulnerability & CVE, Misconfiguration, Logic flaw & Lỗ hổng hoặc điểm yếu đã nhận diện. \\
\hline
Finding & Candidate, Validated, Verified, Reported & Kết luận về một vấn đề an toàn thông tin, có lifecycle bốn bước rõ ràng. \\
\hline
Action & Tool call, Command, Exploit attempt & Bản ghi thao tác đã thực thi, phục vụ audit và replay. \\
\hline
Evidence & stdout, HTTP response, Screenshot, File artifact & Dữ liệu thô hỗ trợ cho finding hoặc state. \\
\hline
\end{tabular}
\end{table}

\paragraph{Các loại cạnh trong AKG}

Quan hệ giữa các node được biểu diễn bằng các cạnh có nhãn, có hướng và trọng số (chi phí, rủi ro OPSEC, mức tin cậy). Mười ba cạnh chính cùng domain/range tương ứng được liệt kê trong Bảng~\ref{tab:akg-edges}.

\begin{table}[H]
\centering
\caption{Các loại cạnh trong AKG với domain và range}
\label{tab:akg-edges}
\small
\renewcommand{\arraystretch}{1.2}
\begin{tabular}{|p{3.2cm}|p{3.4cm}|p{3.4cm}|p{5.5cm}|}
\hline
\textbf{Cạnh} & \textbf{Domain} & \textbf{Range} & \textbf{Ý nghĩa} \\
\hline
\texttt{HOSTS} & Asset & Asset & Host chứa service/port \\
\hline
\texttt{RUNS\_ON} & Asset & Asset & Service chạy trên host \\
\hline
\texttt{EXPOSES} & Asset & Vulnerability & Asset có lỗ hổng \\
\hline
\texttt{EXPLOITABLE\_VIA} & Vulnerability & Action & Lỗ hổng được khai thác qua action \\
\hline
\texttt{YIELDS} & Action & Credential \textbar\ Evidence & Action sinh ra credential hoặc evidence \\
\hline
\texttt{LEADS\_TO} & Action & Finding & Action dẫn tới finding \\
\hline
\texttt{EVIDENCES} & Evidence & Finding & Evidence hỗ trợ finding \\
\hline
\texttt{HAS\_CREDENTIAL} & Identity & Credential & Identity sở hữu credential \\
\hline
\texttt{CAN\_AUTH\_AS} & Credential & Identity & Credential cho phép xác thực thành identity \\
\hline
\texttt{MEMBER\_OF} & Identity & Identity & Thuộc nhóm/vai trò \\
\hline
\texttt{HAS\_PRIVILEGE} & Identity & Asset & Identity có quyền trên asset \\
\hline
\texttt{CONNECTS\_TO} & Asset & Asset & Kết nối mạng giữa các asset \\
\hline
\texttt{PIVOT\_VIA} & Asset & Asset & Pivot qua asset trung gian \\
\hline
\end{tabular}
\end{table}

\paragraph{Ví dụ một đoạn AKG}

\begin{figure}[H]
\centering
\begin{tikzpicture}[
    node distance=10mm and 14mm,
    n/.style={draw, rounded corners, align=center, inner sep=3pt, font=\scriptsize, minimum height=8mm},
    e/.style={-{Latex[length=1.6mm]}, thick, font=\scriptsize}
]
\node[n, fill=red!15] (webapp) {WebApp\\\texttt{/upload}};
\node[n, fill=red!15, right=of webapp] (cve) {CVE-2025-XXXX\\RCE};
\node[n, fill=orange!15, right=of cve] (shell) {Action:\\exec\_rce};
\node[n, fill=blue!15, below=of shell] (cred) {Credential\\svc\_backup};
\node[n, fill=blue!15, left=of cred] (smb) {SMB\\fs01};
\node[n, fill=green!15, left=of smb] (find) {Finding\\Verified};

\draw[e] (webapp) -- node[above]{\texttt{EXPOSES}} (cve);
\draw[e] (cve) -- node[above]{\texttt{EXPLOITABLE\_VIA}} (shell);
\draw[e] (shell) -- node[right]{\texttt{YIELDS}} (cred);
\draw[e] (cred) -- node[above]{\texttt{CAN\_AUTH\_AS}} (smb);
\draw[e] (smb) -- node[above]{\texttt{EVIDENCES}} (find);
\end{tikzpicture}
\caption{Ví dụ một đoạn AKG sau pha Initial Foothold}
\label{fig:akg-example}
\end{figure}

\paragraph{Cơ chế cập nhật và tính nhất quán}

AKG ghi event theo mô hình \textbf{append-only} (bất biến) và state là ảnh chụp tái tạo được từ event stream (versioned). Kết hợp hai nguyên tắc này đảm bảo: auditability (truy vết finding về event gốc), replayability (dựng lại toàn mission để debug), concurrency (ghi song song qua transaction theo node) và recovery (phục hồi sau crash mà không mất dữ liệu).

\paragraph{Ba invariant của AKG}

Để đảm bảo tính đúng đắn của đồ thị trong suốt mission, AKG được bảo vệ bởi ba invariant:

\begin{itemize}
    \item \textbf{Schema invariant:} mọi node và edge phải tuân theo schema đã định nghĩa; mọi edge phải tham chiếu tới hai endpoint tồn tại và đúng kiểu theo schema; các node bắt buộc phải thỏa ràng buộc định danh và provenance tương ứng.
    \item \textbf{Provenance invariant:} mọi finding phải truy vết được về evidence và action đã tạo ra nó. Finding không có evidence bị coi là chưa xác minh.
    \item \textbf{Scope invariant:} mọi node asset phải nằm trong scope đã được phê duyệt; node ngoài scope bị từ chối ngay tại thời điểm ghi. Đây là lớp bảo vệ thứ hai, độc lập với scope guard ở tầng tool executor.
\end{itemize}

\paragraph{Lựa chọn backend}

AKG không phụ thuộc backend cụ thể. Trong giai đoạn thực tập, Neo4j được chọn làm backend chính vì tương thích với Decepticon và hỗ trợ Cypher-native path search với APOC; Memgraph (tương thích Cypher) và graph store nhúng (NetworkX/igraph) là lựa chọn thay thế cho môi trường streaming hoặc lab nhẹ.
% =====================================================================
%  File: sections/3_4_3_mo_hinh_hai_pha.tex
%  Mục đích: Phần 3.4.3 - Mô hình khai thác hai pha
% =====================================================================

\subsubsection{Mô hình khai thác hai pha}
\label{sec:mo-hinh-hai-pha}

Điểm khác biệt quan trọng thứ hai của Raptix so với Decepticon gốc là việc phân tách quy trình khai thác thành hai pha với \emph{chiến lược lan truyền} khác nhau. Sự tách biệt này phản ánh đúng thực tế tác chiến red team: giai đoạn đầu cần \emph{độ phủ} (breadth) để tìm điểm vào, giai đoạn sau cần \emph{độ sâu} (depth) để đạt mục tiêu cuối.

\paragraph{Pha 1 - Initial Foothold}

Agent quét bề mặt tấn công theo chiều ngang (breadth-first, song song hóa cao) để liệt kê asset, service và điểm vào tiềm năng. AKG ghi các node Asset, Service, Vulnerability (candidate), Credential qua cạnh \texttt{EXPOSES}, \texttt{HOSTS}, \texttt{RUNS\_ON}. Pha kết thúc khi có ít nhất một foothold được xác minh và bề mặt đủ để lập kế hoạch pha sau.

\paragraph{Pha 2 - Network Propagation}

Từ foothold, agent chuyển sang depth-first (song song hóa thấp, tuần tự theo attack path) để leo thang đặc quyền, thu credential và di chuyển ngang. AKG ghi node Credential, Identity (đặc quyền được biểu diễn qua cạnh \texttt{HAS\_PRIVILEGE}) qua cạnh \texttt{CAN\_AUTH\_AS}, \texttt{MEMBER\_OF}, \texttt{HAS\_PRIVILEGE}, \texttt{PIVOT\_VIA}. Pha kết thúc khi đạt mục tiêu cuối hoặc hết ngân sách tài nguyên.

\paragraph{Cơ chế chuyển pha}

Điểm cốt lõi của mô hình hai pha là \textbf{cơ chế chuyển pha tường minh dựa trên trạng thái của AKG}. Cụ thể, Raptix duy trì một tập các điều kiện chuyển pha được đánh giá liên tục sau mỗi bước hành động:

\begin{itemize}
    \item \textbf{Điều kiện chuyển từ pha 1 sang pha 2:} AKG chứa ít nhất một node Finding thuộc loại foothold đã đạt trạng thái \texttt{verified}, đồng thời bề mặt tấn công đã được liệt kê tới mức đủ để planner có thể đề xuất ít nhất một attack path khả thi tới mục tiêu cuối.
    \item \textbf{Điều kiện quay lại pha 1:} trong pha 2, nếu agent gặp một subnet mới chưa được liệt kê (thể hiện bằng node Asset không có cạnh \texttt{HOSTS} hoặc \texttt{CONNECTS\_TO} tới các node đã biết), planner có thể tạm thời chuyển sang chế độ breadth-first để mở rộng bề mặt, sau đó quay lại depth-first.
    \item \textbf{Điều kiện kết thúc:} mục tiêu cuối đã đạt được, hoặc ngân sách token/thời gian đã hết, hoặc operator ra lệnh dừng.
\end{itemize}

Mọi chuyển pha đều được ghi lại trong AKG như một event có timestamp, lý do chuyển pha và trạng thái đồ thị tại thời điểm chuyển. Điều này cho phép audit và replay toàn bộ quá trình chuyển pha.

\paragraph{Tận dụng AKG cho attack vector reasoning}

Trong pha 2, AKG đóng vai trò công cụ suy luận chủ động: (i) truy vấn đường đi từ foothold tới mục tiêu bằng weighted shortest-path (các thuộc tính chi phí, rủi ro OPSEC và confidence được chuẩn hóa thành một trọng số tổng hợp $w(e)$ trước khi chạy Dijkstra); (ii) phát hiện credential reuse path qua \texttt{HAS\_CREDENTIAL} và \texttt{CAN\_AUTH\_AS}; (iii) phát hiện privilege escalation path qua \texttt{MEMBER\_OF} và \texttt{HAS\_PRIVILEGE}; (iv) loại trừ nhánh đã thất bại (đánh dấu \texttt{status: failed}); (v) ưu tiên hóa theo chi phí và rủi ro OPSEC.

\begin{figure}[H]
\centering
\begin{tikzpicture}[
    node distance=8mm and 12mm,
    n/.style={draw, rounded corners, align=center, inner sep=3pt, font=\scriptsize, minimum height=8mm},
    e/.style={-{Latex[length=1.6mm]}, thick},
    e2/.style={-{Latex[length=1.6mm]}, dashed, gray}
]
\node[n, fill=orange!20] (foothold) {Foothold\\web01};
\node[n, fill=blue!15, right=of foothold] (cred1) {Cred\\svc\_web};
\node[n, fill=blue!15, right=of cred1] (fs) {Host\\fs01};
\node[n, fill=blue!15, right=of fs] (da) {DA\\dc01};
\node[n, fill=green!20, above=of fs] (goal) {Goal\\Domain Admin};

\draw[e] (foothold) -- node[above]{\texttt{PIVOT\_VIA}} (cred1);
\draw[e] (cred1) -- node[above]{\texttt{CAN\_AUTH\_AS}} (fs);
\draw[e] (fs) -- node[above]{\texttt{HAS\_PRIVILEGE}} (da);
\draw[e] (da) -- (goal);

\node[n, fill=gray!20, below=of fs] (fail) {Đã thử -- thất bại};
\draw[e2] (cred1) -- (fail);
\end{tikzpicture}
\caption{Ví dụ attack path depth-first trong pha Network Propagation; nhánh thất bại được đánh dấu để tránh lặp lại}
\label{fig:attack-path-example}
\end{figure}

\paragraph{So sánh hai pha}

\begin{table}[H]
\centering
\caption{So sánh đặc trưng của hai pha khai thác}
\label{tab:so-sanh-hai-pha}
\small
\renewcommand{\arraystretch}{1.3}
\begin{tabular}{|p{3.2cm}|p{6cm}|p{6cm}|}
\hline
\textbf{Tiêu chí} & \textbf{Initial Foothold} & \textbf{Network Propagation} \\
\hline
Chiến lược & Breadth-first, ưu tiên độ phủ & Depth-first, ưu tiên đường tới mục tiêu \\
\hline
Mục tiêu & Điểm vào đầu tiên + bản đồ bề mặt & Mục tiêu cuối (DA, data, flag) \\
\hline
Song song hóa & Cao, nhiều target cùng lúc & Thấp, phụ thuộc tuần tự theo attack path \\
\hline
Vai trò AKG & Ghi nhận asset, service, vuln & Ghi nhận credential, privilege, pivot \\
\hline
Rủi ro chính & Bỏ sót bề mặt tấn công & Đi sai nhánh, tốn tài nguyên \\
\hline
\end{tabular}
\end{table}

